\chapter{Principio de No-Sobrecarga: Integración sin agregar tiempo}
\label{anexo:principio-no-sobrecarga}

\section{1. Principio de No-Sobrecarga: Integración sin agregar tiempo}

Una preocupación frecuente ante cualquier extensión metodológica es el tiempo adicional que consume. Esta sección demuestra que SocialScrum, implementado correctamente, agrega una sobrecarga mínima: aproximadamente 45 minutos cada dos semanas, equivalente al 0.6\% del tiempo total del Sprint.

El principio fundamental es \textbf{integración, no adición}: los componentes sociales de SocialScrum se incorporan a los eventos Scrum existentes, no crean reuniones nuevas.

\subsection{Sobrecarga temporal por evento}

\begin{table}[!htbp]
    \centering
    \small
    \begin{tabularx}{\textwidth}{|X|c|c|c|c|}
        \hline
        \textbf{Evento} & \textbf{Scrum estándar} & \textbf{SocialScrum} & \textbf{Sobrecarga} (redistribución interna) & \textbf{\% Aumento} \\
        \hline
        Social Sprint Planning
                        & 4 horas
                        & 4 h 15 min
                        & +15 min
                        & +6\%                                                                                                                \\
        \hline
        Social Daily Standup (x10)
                        & 150 min total
                        & 170 min total
                        & +20 min
                        & +13\%                                                                                                               \\
        \hline
        Social Sprint Review
                        & 2 horas
                        & 2 h 10 min
                        & +10 min
                        & +8\%                                                                                                                \\
        \hline
        Social Sprint Retrospective
                        & 1.5 horas
                        & 1.5 horas
                        & +0 min
                        & 0\%                                                                                                                 \\
        \hline
        \textbf{Total por Sprint}
                        & \textbf{10 horas}
                        & \textbf{10 h 45 min}
                        & \textbf{+45 min}
                        & \textbf{+7.5\%}                                                                                                     \\
        \hline
    \end{tabularx}
    \caption{Sobrecarga temporal de SocialScrum por evento}
    \label{tab:sobrecarga-temporal}
\end{table}
\FloatBarrier

\subsubsection{Detalle por evento}

\begin{table}[!htbp]
    \centering
    \small
    \begin{tabularx}{\textwidth}{|l|X|c|}
        \hline
        \textbf{Evento} & \textbf{Cambio en SocialScrum}                                                                                       & \textbf{Tiempo adicional} \\
        \hline
        Social Sprint Planning
                        & Tema 3: Riesgos sociales (Health Score, \textit{community smells} activos e ítems sociales)
                        & +15 min                                                                                                                                          \\
        \hline
        Social Daily Standup (por día)
                        & Pregunta final: “¿Hay algún impedimento social hoy?” (solo identificar, no resolver)
                        & +2--3 min/día                                                                                                                                    \\
        \hline
        Social Sprint Review
                        & Sección dedicada a la evaluación de la dinámica social y la correlación entre Health Score y velocidad
                        & +10 min                                                                                                                                          \\
        \hline
        Social Sprint Retrospective
                        & Redistribución del tiempo: inspección técnica (30 min), inspección social (45 min) y definición de acciones (15 min)
                        & +0 min                                                                                                                                           \\
        \hline
    \end{tabularx}
    \caption{Detalle de cambios por evento}
    \label{tab:detalle-overhead}
\end{table}
\FloatBarrier

\subsection{Sesiones adicionales y presupuesto}

Algunos items sociales requieren sesiones fuera de eventos estándar (mediaciones 1:1, talleres de cohesión). Estas sesiones sí representan tiempo adicional, pero están acotadas:
\begin{itemize}
    \item \textbf{Máximo 20\% de capacidad del equipo} para items sociales.
    \item \textbf{Máximo 2 horas de sesiones adicionales por semana} (fuera de eventos Scrum).
    \item Si se requiere más, es señal de crisis que debe escalarse a CTO.
\end{itemize}

\subsection{Optimizaciones prácticas}

\begin{table}[!htbp]
    \centering
    \small
    \begin{tabularx}{\textwidth}{|l|X|c|}
        \hline
        \textbf{Optimización} & \textbf{Descripción}                                                                                     & \textbf{Ahorro} \\
        \hline
        Daily asíncrono
                              & Reemplazar el Daily síncrono por un formato escrito en Slack o Teams
        (campos: DONE, TODO, BLOCK, SOCIAL)
                              & Elimina 15--18 min/día de reunión                                                                                          \\
        \hline
        Retros alternadas
                              & Sprint impar: énfasis técnico (50/25 min). Sprint par: énfasis social (25/50 min)
                              & Reduce la fatiga de retrospectiva                                                                                          \\
        \hline
        Social Guide rotativo
                              & En equipos de 5--7 personas, rotar el rol cada Sprint.
        Dedicación aproximada: 10\% del tiempo cuando le corresponde
                              & Distribuye la carga entre todos                                                                                            \\
        \hline
        Integración primero
                              & Jerarquía de integración: Daily $\rightarrow$ Planning $\rightarrow$ Review $\rightarrow$ Retrospective.
        Solo crear una sesión adicional si no cabe en los eventos existentes
                              & Minimiza sesiones adicionales                                                                                              \\
        \hline
    \end{tabularx}
    \caption{Optimizaciones para reducir sobrecarga}
    \label{tab:optimizaciones-overhead}
\end{table}
\FloatBarrier

\subsection{Gestión de expectativas}

\subsubsection{Comunicación al equipo}

Mensaje clave: ``SocialScrum agrega 45 min cada 2 semanas a nuestras reuniones Scrum (4.5 min/día). No creamos reuniones nuevas; nos integramos en las existentes. Beneficio esperado: equipos con mejor comunicación entregan más a largo plazo.''
\subsubsection{Comunicación a stakeholders}

Mensaje clave: ``Implementamos SocialScrum para mejorar sostenibilidad del equipo. Verán un Health Score en cada Sprint Review. Sobrecarga: 0.6\% del tiempo. ROI: una renuncia evitada ahorra 3-6 meses de capacitación del reemplazo.''

\subsection{Retorno de inversión (ROI) de SocialScrum}

El retorno de inversión se estima en dólares y pesos colombianos con conversión de \$1 dólar es igual a \$4.000 COP, en este escenario se hace un ejemplo mostrado en las tablas \ref{tab:inversion-anual} y \ref{tab:roi-socialscrum}.

\begin{table}[!htbp]
    \centering
    \small
    \begin{tabularx}{\textwidth}{|l|c|X|c|c|}
        \hline
        \textbf{Inversión (anual)} & \textbf{Horas}                  & \textbf{Detalle} & \textbf{Costo (USD)} & \textbf{Costo (COP)} \\
        \hline
        Sobrecarga de reuniones
                                   & 19.5 h
                                   & 45 min $\times$ 26 Sprints
                                   & \$975
                                   & $\sim$3\,900\,000 COP                                                                            \\
        \hline
        Sesiones adicionales
                                   & 130 h
                                   & $\sim$5 h $\times$ 26 Sprints
                                   & \$6\,500
                                   & $\sim$26\,000\,000 COP                                                                           \\
        \hline
        \textbf{Total}
                                   & \textbf{150 h}
                                   & Costo anual estimado
                                   & \textbf{\$7\,500}
                                   & \textbf{$\sim$30\,000\,000 COP}                                                                  \\
        \hline
    \end{tabularx}
    \caption{Inversión anual estimada en sobrecarga de SocialScrum}
    \label{tab:inversion-anual}
\end{table}
\FloatBarrier

\begin{table}[!htbp]
    \centering
    \small
    \begin{tabularx}{\textwidth}{|l|X|}
        \hline
        \textbf{Retorno (anual)} & \textbf{Valor estimado}                                     \\
        \hline
        1 renuncia evitada
                                 & \$50\,000--\$100\,000 dólares (50--200\% del salario anual) \\
        \hline
        Mejora de velocidad 10\%
                                 & 104 SP adicionales (4 SP por Sprint $\times$ 26 Sprints)    \\
        \hline
        Conflictos resueltos
                                 & Más de 50 horas de fricción evitada                         \\
        \hline
    \end{tabularx}
    \caption{Análisis de ROI de SocialScrum}
    \label{tab:roi-socialscrum}
\end{table}
\FloatBarrier

\textbf{Balance}: ROI positivo si se evita 1 renuncia anual o se mejora velocidad 5\%+.